\chapter{API Specification}
\label{app:api}

This appendix provides the REST API specification for deployment of the feasibility prediction model.

%==============================================================================
\section{Overview}
%==============================================================================

The API exposes three endpoints:
\begin{itemize}
\item \texttt{POST /predict} --- Feasibility prediction from lab/vital input
\item \texttt{GET /model} --- Model metadata and version
\item \texttt{POST /explain} --- SHAP explanation for a prediction
\end{itemize}

Base URL: \texttt{https://api.example.com/v1}

%==============================================================================
\section{Endpoints}
%==============================================================================

\subsection{POST /predict}

\textbf{Request:}
\begin{lstlisting}
{
  "patient_id": "string",
  "features": {
    "creatinine": 1.2,
    "glucose": 120,
    "hemoglobin": 10.5
  }
}
\end{lstlisting}

\textbf{Response:}
\begin{lstlisting}
{
  "patient_id": "string",
  "feasible": true,
  "probability": 0.72,
  "model_version": "1.0"
}
\end{lstlisting}

\subsection{GET /model}

Returns model metadata: version, AUROC, training date, feature list.

\subsection{POST /explain}

Same request body as \texttt{/predict}. Returns SHAP values per feature for the prediction.

%==============================================================================
\section{OpenAPI Summary}
%==============================================================================

\begin{table}[htbp]
\centering
\caption{API Endpoints Summary}
\begin{tabular}{lll}
\toprule
Method & Path & Description \\
\midrule
POST & /predict & Feasibility prediction \\
GET & /model & Model metadata \\
POST & /explain & SHAP explanation \\
\bottomrule
\end{tabular}
\label{tab:api_endpoints}
\end{table}

Authentication: API key in \texttt{X-API-Key} header. Rate limit: 100 req/min.
